% =====================================================================
%  4 Verhalten  (Aufgabenstellung: „Visualisierung eines zentralen
%  Systemablaufs der Anwendung durch ein UML-Sequenzdiagramm“)
%  Tipp: denselben Kernprozess wie im Aktivitätsdiagramm der
%  Spezifikation wählen, nun auf Ebene der Klassen/Objekte.
% =====================================================================
\section{Verhalten}
\TODO{Kurze Hinführung}

\subsection{Zentraler Systemablauf}
\TODO{Ablauf beschreiben (Auslöser, beteiligte Objekte, Nachrichten, Ergebnis)}

\subsection{Sequenzdiagramm}
\TODO{UML-Sequenzdiagramm in docs/diagramme/05\_Sequenzdiagramm/ – Export als docs/diagramme/export/sequenz.pdf}
\begin{figure}[H]
  \caption{UML-Sequenzdiagramm des zentralen Systemablaufs}\label{abb:sequenz}
  \centering
  \grafik{sequenz}
  \quelle{Eigene Darstellung.}
\end{figure}
